\section{Discussion}

\noindent Les résultats expérimentaux montrent que l'architecture
CNN-BiLSTM-Attention proposée est efficace pour détecter les fuites d'eau dans
des séries temporelles multivariées. Le modèle atteint un fort pouvoir
discriminant et offre un compromis équilibré entre précision, rappel et taux de
fausses alertes, ce qui est essentiel pour les systèmes de surveillance réels.

\vspace{0.2cm}

\noindent La comparaison de benchmark met en évidence des différences clés entre les
approches. Bien que Random Forest obtienne une AUC-PR élevée, il présente un
rappel et un score F1 plus faibles au point de fonctionnement sélectionné. La
régression logistique est compétitive grâce à l'ingénierie de caractéristiques,
mais elle ne modélise pas correctement les dépendances temporelles. En revanche,
le modèle proposé offre le meilleur équilibre global, confirmant l'intérêt de
combiner des mécanismes convolutionnels, récurrents et d'attention.

\vspace{0.2cm}

\noindent Un autre aspect important est l'utilisation de la calibration du seuil,
qui permet au modèle de s'adapter à différentes exigences opérationnelles en
contrôlant le compromis entre sensibilité de détection et fausses alertes.
Cette flexibilité est essentielle pour un déploiement pratique dans des systèmes
réels.

\vspace{0.2cm}

\noindent Malgré ces résultats prometteurs, certaines limites subsistent. Le jeu de
données est restreint à une période de surveillance spécifique, ce qui peut
limiter la généralisation. De plus, la comparaison de benchmark pourrait être
étendue à davantage de modèles, et l'interprétabilité de l'architecture de deep
learning demeure limitée.

\vspace{0.2cm}

\noindent Globalement, les résultats confirment que les architectures hybrides de
deep learning sont bien adaptées à la détection de fuites en séries temporelles,
en particulier lorsqu'elles sont combinées à des stratégies robustes de
prétraitement et d'évaluation.